%document class
\documentclass[10pt,oneside]{book}
%%%% Page Info + Commands %%%%%{

%packages
\usepackage{pgfplots}
\pgfplotsset{compat=1.18}
\usepackage{amsmath}
\usepackage{amsfonts}
\usepackage{amsthm,letltxmacro}
\usepackage{amssymb}
\usepackage{enumerate}
\usepackage{enumitem}
\usepackage[dvipsnames]{xcolor}
\usepackage{changepage}
\usepackage{mdframed}

%This will shrink the page
\newcommand{\smallpage}{  \setlength \oddsidemargin{.25in}
            \setlength \textwidth{5.5in}
            \setlength \topmargin{0in}
            \setlength \textheight{9in}}


% Commands for sets of numbers
\newcommand{\Z}{\mathbb{Z}}\newcommand{\R}{\mathbb{R}}\newcommand{\N}{\mathbb{N}}

% gordie spacing and boxing commands
\newcommand{\gspc}{\vspace{0.13cm}} % 0.5 space
\newcommand{\gline}{\gspc\\} % 1.5 space
\newcommand{\gbline}{\gspc\gspc\\} % double space
\newcommand{\gbox}[1]{\fbox{\parbox{\linewidth}{#1}}} % multiline box command
\newcommand{\gmod}[1]{\,(\text{mod}\,#1)}


\newcommand{\gimage}[3]{
\begin{figure}[h]   
    \centering          
    \includegraphics[width=#2\textwidth]{#1}  
    \caption{#3}
    \label{fig:#1}
\end{figure}\\
}

\newcommand{\centerit}[1]{{\begin{center} #1 \end{center}}}
\newcommand{\ul}[1]{\underline{#1}}

%%%%%%%% command for graphics %%%%%%%%%%%%%
\usepackage{fancyhdr}
%}

%%%% Page Setup %%%%%%%
\smallpage\sloppy
\pagestyle{fancy}
\lhead{\large{\textbf{Problem Set 10}}}
%\rfoot{\thepage/\pageref{LastPage} }
\setlength{\headheight}{10pt}

%coloring with color
\newmdenv[
  backgroundcolor=gray!8,
  linewidth=0pt,
  innerleftmargin=0pt,
  innerrightmargin=0pt,
  skipabove=\baselineskip,
  skipbelow=\baselineskip
]{coloredadjust}

% Proof writing
\LetLtxMacro\oldproof\proof
\let\endoldproof\endproof
\renewenvironment{proof}[1][\proofname]
  {\vspace{0.2cm}\begin{adjustwidth}{2em}{2em}
   \oldproof[#1]}
  {\endoldproof
\end{adjustwidth}}


% problem writing
\newenvironment{problem}[1]
  {\vspace{-0.1cm}\begin{coloredadjust}\begin{adjustwidth}{2em}{2em}
   \textit{Problem. }#1}
  {\endoldproof
\end{adjustwidth}\end{coloredadjust}\gspc}

\begin{document}

\newcommand{\cg}[1]{\color{OliveGreen}#1\color{white}}
\newcommand{\cb}[1]{\color{BlueViolet}#1\color{white}}

\subsection*{Section 9.1}

\begin{enumerate}
	\item Solve the following quadratic congruences:
	\begin{enumerate}
		\item [(b)]$3x^2 + 9x + 7 = 0 \gmod{13}$
    \begin{problem}
      We solve this problem by utilities the ``$y^2\equiv d\gmod{p}$'' method from the textbook. First I will do some quick algebra showing why this method works:
      \begin{align*}
        ax^2 + bx + c &= 0\\
        4a(ax^2+bx + c) &= 0\\
        4a^2x^2+4abx + 4ac &= 0\\
        (2ax+b)^2 - b^2 + 4ac &= 0\\
        (2ax+b)^2 &= b^2 - 4ac
      \end{align*}
      Thus, we rewrite our problem in the form:
      \begin{align*}
        y^2 = d \gmod{p}
      \end{align*}
      For $y = (2ax+b)$ and $d = b^2-4ac$.\gline
      Now we can begin with the problem. Consider that in this case:
      \begin{align*}
        y &= 6x + 9\\
        d&= -3
      \end{align*}
      Rewriting gives:
      \begin{align*}
        (6x+9)^2 &\equiv 10 \gmod{13}
      \end{align*}
      Now, we manually solve for $y$-values. We find that $y = 6$ gives $y^2 = 36 \equiv 10 \gmod{13}$. Solving this equation gives with $10 = 2x + b$ gives that $x = 6$. Subtracting $13 - 6$ gives the other solution for $y = 7$.  
    \end{problem}
		\item [(c)]$5x^2 + 6x + 1=0\gmod {23}$
    \begin{problem}
      We'll use the same solution method for this problem too. Consider the following $y$ and $d$:
      \begin{align*}
        y &= 10x+ 6\\
        d &= 16
      \end{align*}
      Trying out various $y$ values gives us that $y = 4$. Solving for $x$ gives that $10x + 6 \equiv 4\gmod{23}$. Thus, $10x\equiv 21\gmod{23}$, and thus using the euclidean:
      \begin{align*}
        23 &= 2(10) + 3\\
        10 &= 3(3) + 1\\
        3 &= 3(1) + 0
      \end{align*}
      With reverse:
      \begin{align*}
        10 - 3(3) &= 1\\
        10-3(23 - 2(10)) &= 1\\
        7(10) - 3(23) &= 1
      \end{align*}
      Gives us that 
      $x = 7\cdot 21 = 147 \equiv 9\gmod{23}$. Thus one of our solutions is $x = 9$. \gline 
      The other solution is given from $10x + 6 = -4$, such that $10x \equiv -10\gmod{23}$. This very simply gives $x = -1 \equiv 22 \gmod{23}$, \gline
      Thus our two solutions are $x =9, x = 22$. 
    \end{problem}
	\end{enumerate}
  \item Prove that the quadratic congruence $6x^2 + 5x + 1 \equiv 0\gmod{p}$ has a solution for every prime $p$. \gline
  I'm going to do this a little bit weird because we can factor this to get the solutions early rather than going through a lot of the legwork. 
  \begin{proof}
    Let $p$ be a prime. We want to show that the quadratic congruence $6x^2 + 5x + 1 \equiv 0\gmod{p}$ has a solution.\gline
    Consider the factorization $(3x+1)(2x+1)$. Thus we have that for solutions, we have that:
    \begin{align*}
      3x \equiv -1 \gmod{p}\\
      2x \equiv -1 \gmod{p}
    \end{align*}
    \centerit{Case: $p = 2$}
    Consider the case where $p = 2$. Thus we have that $3x\equiv -1\gmod{2}$ has a solution for $x = -1$. 
    \centerit{Case: $p =3$}
    Consider the case where $p = 3$. Thus we have that $2x \equiv -1\gmod{3}$ has a solution for $x = 2$. 
    \centerit{Case: $p > 3$}
    Consider the case where $p>3$. Thus, because $\gcd(2, p)=1$ and $\gcd(3,p) = 1$, we have that both equations $3x \equiv -1 \gmod{p}$ and $2x \equiv -1\gmod{p}$ have solutions. \gline
    Thus, the quadratic congruence $6x^2 + 5x + 1 \equiv 0\gmod{p}$ has a solution for every prime $p$, as desired.
  \end{proof}
  \item [12.] If $n>2$ and $\gcd(a,n)=1$, then $a$ is called a quadratic residue of $n$ whenever there exists an integer $x$ such that $x^2\equiv a \gmod{n}$. Prove that if $a$ is a quadratic reside of $n > 2$, then $a^{\phi(n)/2}\equiv 1 \gmod{n}$. 
  \begin{problem}
    Consider via Euler's Theorem that because $a$ is relatively prime to $n$, that $x$ is also relatively prime to $n$. Consider a prime factor $p$ that did divide $x$ and $n$, then $p$ must also divide $x^2 = a$, a contradiction since $\gcd(a,n)= 1$. \gline
    Thus, we have that $x^{\phi(n)}\equiv 1\gmod{n}$, and resultantly, that $(x^{\phi(n)})^2 = a^{\phi(n)/2}\equiv 1 \gmod{n} $. (Note that $\phi(n)$ is even for $n > 2$, so we know that $\phi(n)/2$ is an integer).
  \end{problem}
\end{enumerate}
\subsubsection{Section 9.1}
\begin{enumerate}
  \item [4.] Show that 3 is a quadratic residue of 23, but a nonresidue of 31.
  \begin{problem}
    Consider 3 under Euler's Criterion for mod 23 and 31. $\gcd(3, 23) = \gcd(3, 31) = 1$. We have that:
    \begin{align*}
      3^{(23 - 1)/2}&\equiv 1\gmod{23}\\
      3^{(31-1)/2}&\equiv -1 \gmod{31}
    \end{align*}
    We can show this via the fact that:
    \begin{align*}
      3^{11} &= 3\cdot 3^5 3^5\\
      &\equiv 3\cdot 13\cdot 13 \gmod {23}\\
      &\equiv 3\cdot 8 \gmod{23}\\
      &\equiv 1 \gmod{23}
    \end{align*}
    And for mod 31 we have that:
    \begin{align*}
      3^{15} &= 3^33^63^6\\
      &\equiv 27\cdot 16 \cdot 16\gmod{31}\\
      &\equiv 27\cdot 8 \gmod{31}\\
      &\equiv 30\gmod{31}\\
      &\equiv -1 \gmod {31}
    \end{align*}
  \end{problem}
\end{enumerate}
\subsubsection*{Section 9.2 (wrong number 4)}
\begin{enumerate}
  \item [4.] \begin{enumerate}
    \item Let $p$ be an odd prime. Show that the Diophantine equation $x^2 + py + a = 0\quad \gcd(a,p) = 1$ has a solution if and only if $(-a/p)=1$.  
    \begin{problem}
      We will solve this problem in two directions.
      \centerit{Forwards: $x^2 + py + a = 0\quad \gcd(a,p)=1$}
      First consider the term $py$, which reduces to $0\gmod{p}$. We thus have that:
      \begin{align*}
        x^2 + py + a\equiv x^2 + a \gmod{p}
      \end{align*}
      Simplifying gives that:
      \begin{align*}
        x^2 \equiv -a \gmod{p}
      \end{align*}
      Thus, the equation has a solution if and only if $-a$ is a quadratic residue. 
      \centerit{Backwards: $(-a/p)=1$.}
      We are given that $-a$ is a quadratic residue in $\gmod p$. Thus, let $y$ such that $y^2\equiv -a\gmod{p}$. Let $x = y$ Thus, we have a solution:
      \begin{align*}
        y^2 + py + a &\equiv y^2 + a \gmod{p}\\
        &\equiv -a + a \gmod{p}\\
        &\equiv 0 \gmod{p}
      \end{align*}
      Hence, the diophantine equation $x^2 + py + a = 0$ has a solution if and only if $(-a/p)= 1$. 
    \end{problem}
  \end{enumerate}
  \item [1.] Find the value of the following Legendre symbols:
  \begin{enumerate}
    \item (19/23)
    \begin{problem}
      We use Euler's criterion, as both 19 and 23 are prime and relatively prime to each other.
      \begin{align*}
        19^{(23-1)/2}&=19^{11}\\
        &= (19^3)^3\cdot 19^2\\
        &\equiv (5)^3 \cdot 16\gmod{23}\\
        &\equiv 10 \cdot 16\gmod{23}\\ 
        &\equiv 22\gmod{23}\\
        &\equiv -1 \gmod {23}
      \end{align*}
    \end{problem}
    \item (-23/59)
    \begin{problem}
      We use Euler's Criterion because 23 is relatively prime to 59 (59 is prime).
      \begin{align*}
        (23^{(59-1)/2}) &= 23^{59}\\
        &= -23\cdot(23^2)^14\\
        &\equiv -23\cdot (-2)^{14}\gmod {59}\\
        &\equiv -23 \cdot ((-2)^7)^2\gmod {59}\\
        &\equiv -23 \cdot (49)^2 \gmod{59}\\
        &\equiv 1 \gmod{59}
      \end{align*}
    \end{problem}
    \item (20/31)
    \begin{problem}
      We use Euler's Criterion because 31 is prime. (and thus relatively prime to all other numbers)
      \begin{align*}
        20^{(30-1)/2}&= 20^{15}\\
        &= (20^3)^5\\
        &\equiv (2)^5\gmod{31}\\
        &\equiv 32 \gmod{31}\\
        &\equiv -1 \gmod{31}
      \end{align*}
    \end{problem}
    \pagebreak
    \item (18/43)
    \begin{problem}
      We use Euler's criterion (yet again, because 43 is prime)
      \begin{align*}
        18^{(41-1)/2}&= 18^{20}\\
        &\equiv (18^5)^{4}\gmod{31}\\
        &\equiv 1^4\gmod{31}\\
        &\equiv 1 \gmod{31}
      \end{align*}
    \end{problem}
    \item (-72/131)
    \begin{problem}
      We use Euler's criterion (yet again, yet again because 131 is prime)
      \begin{align*}
        -72^{(131-1)/2} &= (-72)^{65}\\
        &= -1 \cdot (72)(72^4)^{16}\\
        &\equiv -1 \cdot (72) (8^{16})\gmod{131}\\
        &\equiv -1 \cdot (72)(8^8)^2 \gmod {131}\\
        &\equiv -1 \cdot 72(46)^2 \gmod {131}\\
        &\equiv -1 \cdot 72\cdot 20 \gmod{131}\\
        &\equiv -1 \cdot 130 \gmod{131}\\
        &\equiv -1 \cdot -1 \gmod{131}\\
        &\equiv 1 \gmod{131}
      \end{align*}
    \end{problem}
  \end{enumerate}
\end{enumerate}
\end{document}
